%% Included in decks: slidemage-manual, slidemage-manual-copy
%%

% SlideMage manual, chapter "Content and layout": the six block environments,
% each naming itself, and the floating variants.
\setSlideVersion{1-0-0}
\begin{standardframe}{Blocks}
	{\small A block sets a short piece apart and says what it is: each type has
	its own color and icon, so readers see its role at a glance. Each block below
	is named after its environment and fills the width of its column. To give a
	block a width of your own, for instance to put it beside a table, use its floating
	variant: \SMCode{\textbackslash begin\{SMFloatingBlock\}\{5cm\}\{Title\}} (the first argument is the width; also \SMCode{SMFloatingExampleBlock}, \SMCode{SMFloatingInfoBlock}).\par}
	\medskip
	\begin{columns}[t,onlytextwidth]
		\begin{column}{0.31\linewidth}
			\begin{block}{block}
				Neutral: a statement, a formula
			\end{block}
			\begin{definitionblock}{definitionblock}
				A definition, marked $\triangleq$
			\end{definitionblock}
		\end{column}
		\begin{column}{0.31\linewidth}
			\begin{exampleblock}{exampleblock}
				A worked example
			\end{exampleblock}
			\begin{alertblock}{alertblock}
				A caveat or a warning
			\end{alertblock}
		\end{column}
		\begin{column}{0.31\linewidth}
			\begin{infoblock}{infoblock}
				An aside
			\end{infoblock}
			\begin{referenceblock}{referenceblock}
				A reference or citation
			\end{referenceblock}
		\end{column}
	\end{columns}
\end{standardframe}
